% =====================================================================
\section{Introduction}
% =====================================================================

Software systems are rarely static. Over their lifetime, they are
repeatedly \emph{migrated}: adapted internally so they keep working
correctly as the surrounding technology landscape changes. Typical
examples include moving a codebase to a new framework, upgrading a
library to a newer major version, or changing an underlying data
representation for efficiency or safety reasons. These migrations are
necessary maintenance work, but they tend to be tedious, span huge
codebases, and touch code written by many engineers over many years,
which makes them costly and error-prone when done purely by hand
\cite{hora2015developers, kula2018developers}.

Historically, migrations have been automated with techniques ranging
from regular-expression search-and-replace to purpose-built code
transformation tools to fully custom domain-specific tooling. Each
approach trades off precision, recall, and maintenance cost
differently, and none handles the enormous diversity of coding styles
in a codebase the size of Google's particularly well. The emergence
of Large Language Models (LLMs) capable of reading and generating
code opened up a new possibility: rather than hand-writing exhaustive
transformation rules, an LLM can simply be prompted, in natural
language, to make the required edits itself, while much cheaper
heuristics narrow down \emph{where} in the codebase changes might be
needed.

This paper reports on exactly such a system, built and deployed
inside Google to perform a specific, high-stakes migration: changing
identifier fields in Protocol Buffer schemas from 32-bit integers to
64-bit integers, to avoid identifier space exhaustion as systems grow.
The paper's central contribution is a description of the end-to-end
automated pipeline (reference discovery, automatic categorization,
LLM-driven editing, and multi-stage validation) along with an
evaluation across 39 real migrations carried out by three engineers
over a full year. The study offers a large-scale evaluation of
LLM-assisted code migration in a production environment.
